%% kapa: exam=1 round=36 year=2025 period=2 form=A subject=accounting subject_ko=회계학 seq=24 number=64
%% source: Q-Net 감정평가사 기출문제 (cst00309) artlSeq=5231525 file=raw/1008/5231525/2025년 제36회 감정평가사 1차시험 2교시 문제지 원본.pdf
%% conversion: kordoc 4.24.51 -> tools/gen1.py (auto); review=proofread-sonnet; flags=none
\begin{question}{64}
20×1년 영업을 개시한 (주)감평은 상품구매 ￦10당 1포인트를 고객에게 보상하는 고객충성제도를 운영한다. 각 포인트는 기업의 상품을 미래에 구매할 때 ￦1의 할인과 교환할 수 있다. (주)감평은 20×1년 고객에게 현금 ￦100,000에 상품을 판매하고 동시에 10,000포인트를 부여하였다. (주)감평이 부여한 포인트는 별도의 수행의무에 해당하며, 20×1년 판매대가 ￦100,000에는 상품 및 부여된 포인트의 개별 판매가격이 포함되어 있다. (주)감평은 부여된 포인트의 90 \%가 사용될 것으로 예상하고 있고, 1포인트당 개별 공정가치를 ￦0.8으로 추정한다. (주)감평은 예상 사용률을 반영하여 포인트의 개별 판매가격을 측정한 후, 포인트 개별 판매가격을 차감한 잔여금액을 상품의 개별 판매가격으로 배분한다. 20×1년 실제로 6,000포인트가 사용된 경우, (주)감평이 20×1년 인식할 매출(수익)은? (단, 포인트의 소멸시효는 부여일로부터 3년이다.)
\begin{choices}
  \item ￦92,000
  \item ￦92,800
  \item ￦96,400
  \item ￦97,120
  \item ￦97,600
\end{choices}
\end{question}
\answer{5}
